\chapternotnumbered{Introduction}

Given a smooth oriented manifold $X$, let $\mathcal{P}(X)$ be the pseudo-isotopy group of $X$ defined by 
$$\mathcal{P}(X):=\{f\in \Diff(X\times I)|f_{\partial X\times I\cup X\times{0}}=\id\}$$
The pseudo-isotopy group is closely related to the diffeomorphism group of $X$ by $\mathcal{P}(X)\xrightarrow{F} \Diff(X, \partial): f\to F_f=f|_{X\times 1}$, in that case we call that $f$ results in $F_f$.

Hatcher and Wagoner first studied $\pi_0\mathcal{P}(X)$ in their collected work \cite{AST_1973__6__1_0} and it was later studied by Igusa in \cite{K.Igusa}. They found two invariants for the group, namely, $$\Sigma: \pi_0(P)\to \Wh_2(\pi_1X)\text{ and }\Theta: \ker\Sigma\to \Wh_1(\pi_1X; \mathbb{Z}_2\times \pi_2 X)/\chi(K_3\mathbb{Z}[\pi_1 X])$$
where we have $\Wh_1(\pi_1X; \mathbb{Z}_2\times \pi_2 X)=(\mathbb{Z}_2\times \pi_2 M)[\pi_1 M]/(\beta\cdot [1], \alpha\cdot [\sigma]-\alpha^\tau\cdot [\tau\sigma\tau^{-1}], \alpha,\beta\in \mathbb{Z}_2\times \pi_2 M, \tau,\sigma\in \pi_1 M)$.  These two maps are always well-defined whenever $\text{dim}X\geq 4$, and they showed that $\Sigma$ is always surjective when $\text{dim}X\geq 5$, and $\Theta$ is surjective when $\text{dim}X\geq 5$ and bijective when $\text{dim}X\geq 6$. In \cite{singh2022pseudoisotopiesdiffeomorphisms4manifolds} Singh studied the partial images of $\Sigma$ and $\Theta$ in dimension 4 and showed that both $\Sigma$ and $\Theta$ are stably surjective.

The pseudo-isotopy group is closely related to the diffeomorphism group, for there is a natural fiber bundle 
$$\mathcal{J}=\Diff(X\times I,\partial)\to \mathcal{P}(X)\xrightarrow{F} \Diff_{PI}(X,\partial)$$
where $\Diff_{PI}(X,\partial)\subset \Diff(X,\partial) $ denotes all the diffeomorphisms of $X$ which are pseudo-isotopic to identity. Thus we have \emph{induced} Hatcher-Wagoner invariants $$\Sigma: \pi_0\Diff_{PI}(X,\partial)\to \Wh_2(\pi_1 X)/\Sigma(\mathcal{J})$$ $$\Theta: \pi_0F(\ker\Sigma)\to \Wh_1(\pi_1 X; \mathbb{Z}_2\times \pi_2 X)/(\Theta(\mathcal{J}\cap \ker\Sigma)+\chi(K_3\mathbb{Z}[\pi_1 X])) $$ on the diffeomorphism group of $X$, where $\pi_0F(\ker\Sigma)\subset \pi_0\Diff_{PI}(X,\partial)$

Recently, Gabai, Budney, Gay and Hartman constructed the 4-dimensional barbell diffeomorphism (a specific nontrivial element in $\pi_0\Diff(\mathcal{B}_{2,2}^4,\partial)$ following the notations in \cite{Budney_2025}) and studied the implanted barbell diffeomorphisms in a general $X^4$, especially in $S^1\times D^3$ (\cite{budney2021knotted3ballss4}) and $S^4$ (\cite{Gay_2025}). One of the remarkable results is that Budney and Gabai constructed an infinitely generated subset in $\pi_0(\Diff(S^1\times D^3,\partial)/\Diff(D^4,\partial))$, more explicitly, an inclusion $$\bigoplus_{k\geq 4\in \mathbb{N}} \mathbb{Z} \to  \pi_0(\Diff(S^1\times D^3,\partial)/\Diff(D^4,\partial)): 1_k\to \delta_k$$
where $\{\delta_k\}_{k\in \mathbb{N}}$ are specific implanted barbell diffeomorphisms in $S^1\times D^3$ which we will carefully describe later in this paper.

Suppose we have an implanted barbell in a general $X^4$, i.e. $\beta:S^2\times D^2 \natural S^2\times D^2\hookrightarrow X$, and therefore an implanted barbell diffeomorphism in $\Diff(X,\partial)$, it is known that (see \cite[Proposition 2.6]{Budney_2025}) if one of the implanted core spheres is unknotted in $X^4$, then the implanted barbell diffeomorphism is pseudo-isotopic to identity. One natural question is:

\begin{question*}
    How to compute the induced Hatcher-Wagoner invariants for the given half-unkn\-otted implanted barbell diffeomorphism?
\end{question*}

Once we get an $f_\beta\in \pi_0\mathcal{P}$ resulting in that implanted barbell diffeomorphism, $\Sigma(f_\beta), \Theta(f_\beta)$ are representatives of the two induced Hatcher-Wagoner invariants. In this article, we'll construct two pseudo-isotopies $g_\beta,f_\beta\in \pi_0\mathcal{P}$, both resulting in the implanted barbell diffeomorphism with respect to $\beta$. $g_\beta$ has a Cerf diagram containing a single eye of (1,2)-handle pair and $f_\beta$ has a Cerf diagram containing a single eye of (2,3)-handle pair. The main result will give an explicit formula for the Hatcher-Wagoner invariants of $f_\beta$, which is stated as follows:

\begin{theorem}
    For a half-unknotted implanted barbell $\beta=(R_0, S, \gamma)$ with $S=\partial \beta_0^\bullet$ where $\beta_0^\bullet: D^3\hookrightarrow M$, by finger-pushing $R_0$ along the arc, we can make $\gamma$ short enough such that $\text{int}(\beta_0^\bullet)\cap \gamma =\emptyset$. Now suppose that $\beta_0^\bullet \cap R_0=\bigsqcup _{i=1}^k S_i^1$. Choose $p_i\in S_i^1$ and let $*_0=\gamma\cap S$ be the base point. For each $i$, find a path $\delta^B_i\subset \beta=(R_0, S,\gamma)$ which is a path from $*_0$ to $p_i\in S_i^1$. Also, for each $i$, $S_i^1$ divides $R_0$ into two embedded disks $D_i$ and $D_i'$, where $D_i'$ is the one connected to the arc $\gamma$. Let $D^B_i=D_i$ (see \Cref{fig:general-barbell} for an illustration). Then there is a pseudo-isotopy $f_\beta\in \pi_0\mathcal{P}$ resulting in that implanted barbell diffeomorphism with the Cerf diagram of $f_\beta$ being a single eye of (2,3)-handle pair such that:

    $$\Sigma(f_\beta)=0, \Theta(f_\beta)=\sum_{i=1,...,k} (0, [D^B_i]^{\delta_i^B})\cdot  [\delta_i^B] $$
    Here we identify $\pi_i(M,*_0)$ with $\pi_i(M, \beta_0^\bullet)$ so that $[D_i]\in \pi_2(M,\beta_0^\bullet)$ and $\delta_i^B\in \pi_1(M,\beta_0^\bullet)$. 

    \begin{figure}[!ht]
        \includegraphics[width=0.45\textwidth]{general_barbell_3.png}
        \caption{A general barbell and data needed to compute Hatcher-Wagoner invariants}
        \label{fig:general-barbell}
    \end{figure}
\end{theorem}

\begin{remark}
    In particular, when the given barbell $\beta=(R_0,S,\gamma)$ with $S=\partial\beta_0^\bullet$ and $\gamma:[0,1]\to M,\gamma(0)=*_0=\gamma\cap S$ satisfies $R_0\cap \beta_0^\bullet=\emptyset$, suppose $\gamma|_{(0,1)}\cap \beta_0^\bullet=\{p_i\}_{i=1}^k$ with $p_i=\gamma(t_i)$, then the $f_\beta\in \pi_0\mathcal{P}$ we constructed in the theorem satisfies:$$\Sigma(f_\beta)=0, \Theta(f_\beta)=(0,[R_0^\gamma])\sum_{i=1}^k [\gamma_i]$$where $\gamma_i:=\gamma|_{[0,t_i]}\in \pi_1(M,\beta_0^\bullet)$ and $[R_0^\gamma]\in \pi_2(M, \beta_0^\bullet)$ is obtained by pulling $R_0$ back to $*_0$ along $\gamma$.
\end{remark}

\begin{remark}
    Along the way we also compute the Hatcher-Wagoner invariants for $g_\beta$: $\Sigma(g_\beta)=0, \Theta(g_\beta)=-\bar{\Theta}(f_\beta)$, where $\bar{\cdot}$ is an involution on $\Wh_1(\pi_1, \mathbb{Z}_2\times \pi_2)$ which is described in both \cite{AST_1973__6__1_0} and \cite[section 9.1]{singh2022pseudoisotopiesdiffeomorphisms4manifolds}. We review the involution here:
    \[
      \bar{\cdot}: (n,\sigma)\cdot[\gamma]\to (n+w_2^X(\sigma), -w_1^X(\gamma)\sigma^{\gamma^{-1}})\cdot [\gamma^{-1}],
    \]
    where $w_2^X(\sigma)\in \{0,1\}$ is the second Stiefel-Whitney class of the normal bundle of $\sigma$ in $X$, and $w_1^X(\gamma)\in \{-1,1\}$ is the first Stiefel-Whitney class of the normal bundle of $\gamma$ in $X$.
\end{remark}

\begin{remark}
    In particular, all $\{f_{\delta_k}\}_{k\in \mathbb{N}}$ in $\mathcal{P}(S^1\times D^3)$ lie in the kernel of $\Sigma$ and $\Theta$, this answers a question raised by Powell in \cite[Question 12.6]{powellmapping}. But $\{\delta_k, k\geq 4\}$ are nontrivial diffeomorphisms of $S^1\times D^3$, so $[f_{\delta_k}, k\geq 4]\neq 0\in \pi_0\mathcal{P}(S^1\times D^3)$, which equivalently means that $\{f_{\delta_k},k\geq 4\}$ are nontrivial pseudo-isotopies of $S^1\times D^3$ which can not be detected by Hatcher-Wagoner invariants.
\end{remark}

We generalize the calculation to \emph{half-unknotted immersed barbell diffeomorphisms} which will be defined in \Cref{chap:generalizations-to-immersed-barbell-diffeomorphisms}. The result is completely the same:

\begin{theorem}
    For a half-unknotted immersed barbell $\beta=(R_0,S,\gamma)$ with $\beta_0^\bullet: D^3\hookrightarrow M, S=\partial\beta_0^\bullet $, perturb self-intersections of $R_0$ away from $\beta_0^\bullet$. By finger-pushing $R_0$ along arc, we can make $\gamma$ short enough such that $\text{int}(\beta_0^\bullet)\cap \gamma =\emptyset$. Now suppose that $\beta_0^\bullet \cap R_0=\bigsqcup _{i=1}^k S_i^1$. Choose $p_i\in S_i^1$ and let $*_0=\gamma\cap S$ be the base point. For each $i$, find a path $\delta^B_i\subset \beta=(R_0, S,\gamma)$ which is a path from $*_0$ to $p_i\in S_i^1$. Also, for each $i$, $S_i^1$ divides $R_0$ into two embedded disks $D_i$ and $D_i'$, where $D_i'$ is the one connected to the arc $\gamma$. Let $D^B_i=D_i$. Then the $f_\beta\in \pi_0\mathcal{P}$ we constructed which results in the immersed barbell diffeomorphism w.r.t. $\beta$ satisfies:
    
    $$\Theta(f_\beta)=\sum_{i=1,...,k} (0, [D^B_i]^{\delta_i^B})\cdot  [\delta_i^B]$$
    Here we identify $\pi_i(M,*_0)$ with $\pi_i(M, \beta_0^\bullet)$ so that $[D_i]\in \pi_2(M,\beta_0^\bullet)$ and $\delta_i^B\in \pi_1(M,\beta_0^\bullet)$.    
\end{theorem}

From the main theorem we deduce some corollaries: 
\begin{corollary}
    For any $\sigma\in \pi_2 M$, $\forall \alpha\in \pi_1 M$, there is a half-unknotted immersed barbell $\beta=(R, S, \gamma)$ and $f_\beta\in \ker\Sigma\subset \pi_0\mathcal{P}$ with its Cerf diagram being a single eye of (2,3)-handle pair resulting in the immersed barbell diffeomorphism w.r.t. $\beta$ such that $\Theta(f_\beta)=(0,\sigma)\cdot [\alpha]$.
\end{corollary}

In particular, using results deduced by Singh in \cite{singh2022pseudoisotopiesdiffeomorphisms4manifolds} where Singh showed that for $M=(X_1\# X_2)\times I$ with $X_i$ a closed, oriented, aspherical 3-manifold, there's $K\subset \pi_0\Diff_{PI}(M,\partial)$ and a surjection $K\to \bigoplus_{i\in \mathbb{N}} \mathbb{Z}$, we further show that we can let
\[
  K=\langle\text{implanted half-unknotted barbell diffeomorphisms}\rangle,
\]
that is:

\begin{corollary}
    Let $M=(X_1\# X_2)\times I$ with $X_i$ closed, orientable, aspherical 3-manifold. In this case all $\sigma\in \pi_2 M=\mathbb{Z}[\pi_1 X_1*\pi_1 X_2]$ can be realized by embedded $S^2$ with $w_2^M(\sigma)=0$. Then $$\langle\text{implanted half-unknotted barbell diffeomorphisms}\rangle\subset \pi_0\Diff_{PI}(M,\partial)$$ is infinitely generated and of infinite $\mathbb{Z}$-rank. 
\end{corollary}

The paper is organized as follows: In \Cref{chap:cerf-theory-and-parameterised-handle-constructions-version} we recall the Cerf theory and parametrised handle constructions version of it, especially we explain how a loop of handle constructions of $X\times I$ can result in a pseudo-isotopy $f\in \pi_0\mathcal{P}$ of $X$, and how to see the resulting $F_f\in\Diff(X,\partial)$ directly from the loop of handle constructions. In \Cref{chap:twin-twists-and-relation-with-barbell-diffeomorphisms} we recall the general notion of Montesino twins in $X^4$ and the resulting twin twists in $\Diff_0(X,\partial)$ which were studied for $X=S^4$ in \cite{Gay_2025}. Then we show that if one of the spheres in the Montesino twins is unknotted, then $\exists f\in \mathcal{P}(X)$ resulting in that twin twist with Cerf diagram being a loop of a single cancelling (1,2)-handle pair. Then we show that any implanted barbell diffeomorphism is a twin twist, so together with the above result, for a specific implanted barbell $\beta$, we find $g_\beta\in \mathcal{P}(X)$ with Cerf diagram being a loop of a single cancelling (1,2)-handle pair resulting in that implanted barbell diffeomorphism. In \Cref{chap:changing-from-12-handle-pair-to-23-handle-pair} we follow the essential lemma \cite[Lemma 17]{Gay_2025} to develop a technique of \emph{one parameter version of 0-framed and dotted replacement} to change the pseudo-isotopy $g_\beta\in \pi_0\mathcal{P}$ with Cerf diagram being the above loop of cancelling (1,2)-handle pair to $f_\beta\in \pi_0\mathcal{P}$ with Cerf diagram being a loop of cancelling (2,3)-handle pair. Therefore we get a loop of cancelling (2,3)-handle pair resulting in implanted barbell diffeomorphism of $\beta$. In \Cref{chap:computation-for-the-second-hatcher-wagoner-invariant} we use the technology in \Cref{chap:changing-from-12-handle-pair-to-23-handle-pair} to complete the computations of second Hatcher-Wagoner invariant for a half-unknotted implanted barbell diffeomorphism. In \Cref{chap:generalizations-to-immersed-barbell-diffeomorphisms} we generalize to \emph{immersed barbell diffeomorphisms} and compute the induced Hatcher-Wagoner invariants for a half-unknotted immersed barbell diffeomorphism. As corollaries, we deduce some realization results on what induced Hatcher-Wagoner invariants can be realized by immersed half-unknotted barbell diffeomorphisms.

\section*{Acknowledgements}
The author would like to thank Jean Cerf, Allen Hatcher and John Wagoner for developing the excellent theory about pseudo-isotopy. She also thanks David Gabai, David Gay, Daniel Hartman, Ryan Budney and Olivier Singh for their remarkable works on diffeomorphisms and pseudo-isotopies of 4-manifolds. She also thanks Daniel Hartman and Yicheng Yang for useful conversations along the way. She also thanks Yifei Fan for some of the illustrations in this paper. Most especially, She thanks Jianfeng Lin for being her mentor, making many inspiring discussions,  reading the first version of this paper and giving useful advice. 

The work was carried out during the author's fourth year in Qiuzhen College, Tsinghua University, and thus she expresses her deep gratitude to the institution.
